\subsubsubsubsection{Software Interfaces}

Сайт функционирует посредством использования следующего программного обеспечения:

\begin{enumerate}
    \item \textbf{Веб-сервер}
    \begin{itemize}
        \item \textbf{Наименование:} Apache HTTP Server / Nginx.
        \item \textbf{Назначение:} обработка и перенаправление входящих HTTP/HTTPS-запросов от пользователей, обслуживание статических файлов.
        \item \textbf{Версия:} актуальная стабильная версия на момент разработки.
        \item \textbf{Источник:} открытое программное обеспечение.
    \end{itemize}
    \item \textbf{Система управления базами данных (СУБД)}
    \begin{itemize}
        \item \textbf{Наименование:} MySQL или PostgreSQL.
        \item \textbf{Назначение:} хранение всех данных Сайта и управление доступом к ним.
        \item \textbf{Версия:} актуальная стабильная версия на момент разработки.
        \item \textbf{Источник:} открытое программное обеспечение.
    \end{itemize}
    \item \textbf{Сервер приложений}
    \begin{itemize}
        \item \textbf{Наименование:} среда выполения языка программирования, на котором реализована серверная часть Сайта. Конкретный язык определяется разработчиком и согласовывается с заказчиком.
        \item \textbf{Назначение:} выполнение серверной логики Сайта.
        \item \textbf{Версия:} актуальная стабильная версия на момент разработки.
        \item \textbf{Источник:} открытое программное обеспечение.
    \end{itemize}
    \item \textbf{Операционная система сервера}
    \begin{itemize}
        \item \textbf{Наименование:} Linux. Конкретный дистрибутив определяется разработчиком и согласовывается с заказчиком.
        \item \textbf{Назначение:} обеспечение среды выполнения для работы необходимого программного обспечения.
        \item \textbf{Версия:} актуальная стабильная версия на момент разработки.
        \item \textbf{Источник:} открытое программное обеспечение.
    \end{itemize}
    \item \textbf{Клиентское программное обеспечение}
    \begin{itemize}
        \item \textbf{Наименование:} современный веб-браузер, поддерживающий HTML5, CSS3 и JavaScript. Пример: Google Chrome, Mozilla Firefox, Safari, Яндекс.Браузер.
        \item \textbf{Назначение:} отображение пользовательского интерфейса Сайта и обеспечение взаимодействия с пользователем.
        \item \textbf{Версия и Источник:} зависят от конкретного пользователя, поддерживаются актуальный версии браузеров.
    \end{itemize}
\end{enumerate}